\exercisesection{3}

\begin{enumerate}
\def\labelenumi{\arabic{enumi}.}
\item
  For an integral domain to be factorial, it is necessary and sufficient that there exist a mapping \(x \mapsto s(x)\) of A - \{0\} to N such that \(s(xy) = s(x) + s(y)\) , the relation \(s(x) = 0\) implies that x is invertible in A and finally, for any two elements x, y of A - \{0\}, neither of which divides the other, there exist elements a, b, z, t of A - \{0\} such that \(ax + by = zt\) , \(s(z) < s(x)\) , t being prime to x and to y. (If this condition is fulfilled, show first that every non-zero element of A is a product of extremal elements and then that, for every extremal element p, Ap is a prime ideal; conclude with the aid of Theorem 1 (d)).
\item
  \begin{enumerate}
  \def\labelenumii{(\alph{enumii})}
  \tightlist
  \item
    Let A be an integral domain, the union of a right directed family \((\mathbf{A}_{\lambda})\) of subrings. Suppose that each of the \(A_{\lambda}\) is factorial and that, if \(\lambda \leqslant \mu\) , every extremal element of \(A_{\lambda}\) is extremal in \(A_{\mu}\) . Show that A is a factorial domain whose set of extremal elements is the union of the sets of extremal elements of each of the \(A_{\lambda}\) (cf.~§ 2, Exercise 19 (e)).
  \end{enumerate}
\end{enumerate}

\begin{enumerate}
\def\labelenumi{(\alph{enumi})}
\setcounter{enumi}{1}
\item
  Deduce from (a) that for every factorial domain \(\mathbf{A}\) the polynomial domain \(\mathbf{A}[\mathbf{X}_{\lambda}]_{\lambda \in \mathbf{L}}\) in any family of indeterminates is factorial.
\item
  Let A be a factorial domain such that every domain of formal power series \(A[[X_{1},\ldots,X_{n}]]\) in a finite number of indeterminates is factorial. Show that every domain of formal power series \(A[[X_{\lambda}]]_{\lambda\in L}\) in any family of indeterminates (Algebra, Chapter IV, §5, Exercise 1) is factorial (same method).
\end{enumerate}

\begin{enumerate}
\def\labelenumi{\arabic{enumi}.}
\setcounter{enumi}{2}
\tightlist
\item
  \begin{enumerate}
  \def\labelenumii{(\alph{enumii})}
  \tightlist
  \item
    Let \(A = \bigoplus_{n \geqslant 0} A_n\) be a graded algebra with positive degrees over a field k; suppose that \(A_0 = k\) and that A is a Krull domain. Show that for A to be factorial, it is necessary and sufficient that every graded prime ideal p of A of height l be of the form p = Aa, where a is a homogeneous element (use Exercise 16 of §1). Show that every homogeneous element \(\neq 0\) of A is a product of homogeneous extremal elements.
  \end{enumerate}
\end{enumerate}

\begin{enumerate}
\def\labelenumi{(\alph{enumi})}
\setcounter{enumi}{1}
\tightlist
\item
  Let \(\mathbf{A}\) be a graded \(k\) -algebra satisfying the conditions of (a). Let \(k'\) be an extension of \(k\) and suppose that \(\mathbf{A} \otimes_k k'\) is a factorial domain; show that \(\mathbf{A}\) is factorial. (If a graded ideal \(\mathfrak{a}\) of \(\mathbf{A}\) is such that \(\mathfrak{a} \otimes_k k'\) is principal in \(\mathbf{A} \otimes_k k'\) , show that \(\mathfrak{a}\) is principal; then use (a).)
\end{enumerate}

\begin{enumerate}
\def\labelenumi{\arabic{enumi}.}
\setcounter{enumi}{3}
\tightlist
\item
  \begin{enumerate}
  \def\labelenumii{(\alph{enumii})}
  \tightlist
  \item
    Show that the ring \(\mathbf{A} = \mathbf{Q}[\mathbf{X},\mathbf{Y}] / p\) , where \(p\) is the principal ideal generated by \(\mathbf{X}^2 +\mathbf{Y}^2 -1\) in \(\mathbf{Q}[\mathbf{X},\mathbf{Y}]\) , is a non-factorial Krull domain (if \(x\) is the image of \(\mathbf{X}\) in \(\mathbf{A}\) , show that \(x\) is an extremal element of \(\mathbf{A}\) but that \(\mathbf{A}x\) is not a maximal ideal of \(\mathbf{A}\) ).
  \end{enumerate}
\end{enumerate}

\begin{enumerate}
\def\labelenumi{(\alph{enumi})}
\setcounter{enumi}{1}
\tightlist
\item
  Show that the ring \(\mathbf{A} \otimes_{\mathbf{Q}} \mathbf{Q}(i)\) is factorial (prove that this ring is isomorphic to the quotient of \(\mathbf{Q}(i) [X, Y]\) by the principal ideal ( \(XY - 1\) ); compare with Exercise 3 (b)).
\end{enumerate}

¶ 5. (a) Let k be a Noetherian factorial domain and B the polynomial ring \(k[X_{1},\ldots,X_{n}]\) where \(n\geqslant3\) ; let \(g_{i}\;(0\leqslant i\leqslant r)\) be elements of \(k[X_{3},\ldots,X_{n}]\)

where \(g_0\) is extremal. Writing \(g = \mathbf{X}_1\mathbf{X}_2 - \sum_{i=0} g_i\mathbf{X}_1^i\) , consider the quotient ring \(\mathbf{A} = \mathbf{B}/g\mathbf{B}\) . Show that \(\mathbf{A}\) is factorial. (Let \(\mathbf{S}\) be the multiplicative subset of \(\mathbf{A}\) generated by 1 and the image of \(\mathbf{X}_1\) in \(\mathbf{A}\) ; apply Proposition 3 of no. 4 to \(\mathbf{S}^{-1}\mathbf{A}\) .)

\begin{enumerate}
\def\labelenumi{(\alph{enumi})}
\setcounter{enumi}{1}
\item
  Let \(k\) be a field of characteristic \(\neq 2\) and \(F\) a homogeneous polynomial of the second degree in \(k[\mathbf{X}_1, \ldots, \mathbf{X}_n]\) where \(n \geqslant 5\) , such that the corresponding polynomial function on \(k^n\) is a non-degenerate quadratic form. Show that the domain \(k[\mathbf{X}_1, \ldots, \mathbf{X}_n] / (\mathbf{F})\) is factorial. (Prove first that a homogeneous polynomial \(G\) of the second degree in \(k[\mathbf{X}_1, \ldots, \mathbf{X}_n]\) , such that the corresponding polynomial function is a non-degenerate quadratic form, is extremal for \(n \geqslant 3\) . Then prove the proposition when \(k\) is algebraically closed, using (a); finally pass to the general case with the aid of Exercise 3 (b).)
\item
  If \(F = X_{1}X_{2} - X_{3}X_{4}\) , show that the domain \(k[X_{1}, X_{2}, X_{3}, X_{4}]/(F)\) is not factorial. (Show that the images of the \(X_{i}\) in this ring are extremal elements.)
\end{enumerate}

¶ 6. (a) Let K be an algebraically closed field of characteristic ≠2. Determine the graded prime ideals of the ring

\[
\mathrm{K} [ \mathrm{X}, \mathrm{Y}, \mathrm{Z} ] / (\mathrm{X} ^ {2} + \mathrm{Y} ^ {2} + \mathrm{Z} ^ {2}).
\]

\begin{enumerate}
\def\labelenumi{(\alph{enumi})}
\setcounter{enumi}{1}
\item
  Let \(k\) be an ordered field, \(a, b, c\) elements \(>0\) of \(k\) and \(A\) the ring \(k[X, Y, Z] / (aX^2 + bY^2 + cZ^2)\) . Show that \(A\) is a factorial domain. (Reduce it to the case where \(k\) is a maximal ordered field, using Exercise 3 (b)); then prove, using (a), that every graded prime ideal of \(A\) of height 1 is principal and apply Exercise 3 (a).)
\item
  Let \(k\) be an ordered field, \(a, b\) elements \(> 0\) of \(k\) and \(B\) the ring
\end{enumerate}

\[
k [ \mathrm{X}, \mathrm{Y} ] / (a \mathrm{X} ^ {2} + b \mathrm{Y} ^ {2} + 1).
\]

Show that B is a factorial domain (use (b) and Exercise 17 (a) of § 1.)

\begin{enumerate}
\def\labelenumi{(\alph{enumi})}
\setcounter{enumi}{3}
\tightlist
\item
  Show that the domain \(\mathbf{C} = \mathbf{Q}[\mathrm{X},\mathrm{Y}] / (\mathrm{X}^2 +2\mathrm{Y}^2 +1)\) is factorial, but that the domain \(\mathbf{C}\otimes_{\mathbf{Q}}\mathbf{Q}(i)\) (where \(i^2 = -1\) ) is not factorial.
\end{enumerate}

¶ 7. Let K be a Noetherian factorial domain and F an extremal element of the polynomial ring \(K[X_{1},\ldots,X_{n}]\) ; suppose that, when a weight \(q(i)>0\) ( \(1\leqslant i\leqslant n\) ) is attributed to each \(X_{i}\) , F is isobaric and of weight q\textgreater0. Let A be the domain generated, in an algebraic closure \(\Omega\) of the field of fractions of \(K(X_{1},\ldots,X_{n})\) , by K, the \(X_{i}\) ( \(1\leqslant i\leqslant n\) ) and a root z of the polynomial \(Z^{c}-F\) , where c is an integer prime to q. Show that A is factorial in the two following cases:

\begin{enumerate}
\def\labelenumi{(\arabic{enumi})}
\item
  \(c \equiv 1 \pmod{q}\) ;
\item
  every finitely generated projective K-module is free (which holds for example when K is a field or a principal ideal domain or a local ring). (In the first case, consider the ring of fractions A{[}1/z{]}; show that it is a factorial domain and apply Proposition 3 of no. 4. In the second case, consider an integer d such that \(cd \equiv 1 \pmod{q}\) and let \(z' \in \Omega\) such that \(z = z'^{d}\) ; the domain B = A{[}z'{]} is factorial by virtue of the first case and is a free A-module. Consider B as a graded ring taking \(z'\) to be of weight q, each of the \(X_{i}\) of weight \(cdq(i)\) and reduce it to proving that for two homogeneous elements u, v of A the ideal Au ∩ Av is principal, using Exercise 16 of §1; consider finally the ideal Bu ∩ Bv in B and use Chapter I, §3, no. 6, Proposition 12.) In particular, if K is a field and a, b, c are three integers \textgreater0 which are relatively prime in pairs, the domain
\end{enumerate}

\[
\mathbf {A} = \mathrm{K} [ \mathrm{X}, \mathrm{Y}, \mathrm{Z} ] / (\mathrm{Z} ^ {a} - \mathrm{X} ^ {b} - \mathrm{Y} ^ {c})
\]

is factorial.

¶ 8. Let A be an integral domain and x, y, z three non-zero elements of A, where x is extremal and Ax ∩ Ay = Axy. Let S be the multiplicative set of the \(x^{n}\) ( \(n \geqslant 0\) ) and let B = S \(^{-1}\) A; consider the domains of formal power series A{[}{[}T{]}{]} and B{[}{[}T{]}{]}.

\begin{enumerate}
\def\labelenumi{(\alph{enumi})}
\tightlist
\item
  Let \(i, j, k\) be three integers \(\geqslant 0\) such that \(ijk - ij - jk - ki \geqslant 0\) and \(z^i \in Ax^j + Ay^k\) . Consider in A{[}{[}T{]}{]} the element \(v = xy - z^{i-1}\mathrm{T}\) . Show that there exists an integer \(t > 0\) and a series
\end{enumerate}

\[
v ^ {\prime} = y ^ {t} x ^ {- 1} + b _ {1} x ^ {- 2} \mathrm{T} + \dots + b _ {n - 1} x ^ {- n} \mathrm{T} ^ {n - 1} + \dots
\]

in B{[}{[}T{]}{]} such that \(vv' \in A[[T]]\) (determine the \(b_{n}\) inductively, proceeding in N by intervals of length ij; in the interior of each interval, take

\[
b _ {n + 1} = b _ {n} z ^ {i - 1} x ^ {- 1};
\]

at the end of each interval, use the inequality \(ijk + ij - jk - ki \geqslant 0\) ).

\begin{enumerate}
\def\labelenumi{(\alph{enumi})}
\setcounter{enumi}{1}
\item
  Suppose that \(z^{i-1} \notin Ax + Ay\) . Show that there exists in the ring A{[}{[}T{]}{]} no formal power series of constant term \(y^k\) ( \(k\) an integer \(>0\) ) and which is an element associated with v in B{[}{[}T{]}{]} (calculate the coefficient of T in the product of v and an invertible element of B{[}{[}T{]}{]}).
\item
  Deduce from (a), (b) and Exercise 7 an example of a factorial domain A such that the domain of formal power series A{[}{[}T{]}{]} is not factorial. (In the above notation and assuming that the conditions of (a) and (b) on x, y, z, i, j, k are fulfilled, show that \(vv'\) cannot be a product of extremal factors \(u_{h}\) ( \(1 \leqslant h \leqslant r\) ) in A{[}{[}T{]}{]}; observe that the \(u_{h}\) are formal power series whose constant terms are powers of y. Then consider the ring C = S'\^{}\{-1\}A{[}{[}T{]}{]}, where S' consists of the series whose constant term is in S; B{[}{[}T{]}{]} is the completion of the Zariski ring C; show that \(v' \in C\) and that v and the \(u_{h}\) are extremal in C; obtain finally a contradiction with (b).)
\item
  Deduce from (c) an example of a Noetherian factorial local integral domain A whose completion \(\hat{A}\) is a non-factorial Krull domain.
\end{enumerate}

¶ 9. (a) Let A be a Noetherian domain such that, for every maximal ideal m of A, \(A_{m}\) is a discrete valuation ring. If B is the ring of formal power series \(A[[X_{1},\ldots,X_{n}]]\) , show that, for every maximal ideal n of B, the domain \(B_{n}\) is factorial (consider its completion, using Proposition 8 of no. 9). Deduce that every divisorial ideal of B is a projective B-module.

\begin{enumerate}
\def\labelenumi{(\alph{enumi})}
\setcounter{enumi}{1}
\item
  Let C be a Noetherian ring such that every finitely generated projective C-module is free; show that the ring of formal power series C{[}{[}X{]}{]} has the same property (cf.~Chapter II, § 3, no. 2, Proposition 5).
\item
  Deduce from (a) and (b) that, if A is a principal ideal domain, the domain of formal power series \(A[[X_{1},\ldots,X_{n}]]\) is factorial.
\end{enumerate}

\begin{enumerate}
\def\labelenumi{\arabic{enumi}.}
\setcounter{enumi}{9}
\tightlist
\item
  \begin{enumerate}
  \def\labelenumii{(\alph{enumii})}
  \tightlist
  \item
    A Prüferian (§ 2, Exercise 12) factorial domain is a principal ideal domain.
  \end{enumerate}
\end{enumerate}

\begin{enumerate}
\def\labelenumi{(\alph{enumi})}
\setcounter{enumi}{1}
\tightlist
\item
  A pseudo-Bezoutian (§ 1, Exercise 21) Krull domain is factorial.
\end{enumerate}

\begin{enumerate}
\def\labelenumi{\arabic{enumi}.}
\setcounter{enumi}{10}
\item
  Let K be a field and A the polynomial domain \(K[X,Y]\) , which is factorial; if L is the field \(K(X^{2},Y/X) \subset K(X,Y)\) , show that the domain \(A \cap L\) is not factorial.
\item
  Show Proposition 5 of no. 8 using Chapter III, § 2, no. 8, Corollary 3 to Theorem 1.
\item
  Extend the Corollary to Proposition 7 of no. 8 to the case where the complete Hausdorff local ring A is not an integral domain (use Algebra, Chapter VIII, §6, Exercise 6 (b)).
\item
  Let A be a complete Noetherian local ring whose residue field is of characteristic p \textgreater{} 0. In the ring of formal power series A{[}{[}T{]}{]} consider the elements \(\omega_{n} = (1 - T)^{p^{n}}\) and \(\gamma_{n} = 1 - \omega_{n}\) for every integer n \textgreater{} 0. Show that \(\gamma_{n}\) is, to within a sign, a distinguished polynomial (no. 8); deduce that \(A_{n} = A[[T]]/(\gamma_{n})\) is identified with the algebra over A of the group
\end{enumerate}

\(G_{n} = Z/p^{n}Z\) . Show that the intersection of the principal ideals \((\gamma_{n})\) is reduced to 0; deduce that A{[}{[}T{]}{]} is identified with the inverse limit \(\lim A_{n}\) .

¶ 15. Let K be a field which is complete with respect to the valuation v, A the ring of the valuation, k its residue field and P a polynomial in A{[}X₁, \ldots, Xₙ{]} of total degree d with the following property: there exists an algebraic extension K' of K such that in K'{[}X₁, \ldots, Xₙ{]} P is a product of polynomials of total degree 1. Suppose further that there exist two polynomials Q, R in A{[}X₁, \ldots, Xₙ{]} such that Q is of total degree s and contains a monomial \(aX_1^{s}\) where φ(a) ≠ 0 (φ denoting the canonical homomorphism A → k), R is of degree ≤d - s and \(\overline{\mathbf{P}} = \overline{\mathbf{Q}}\) . \(\overline{\mathbf{R}}\) (notation of Chapter III, § 4). Show that there then exist in A{[}X₁, \ldots, Xₙ{]} two polynomials Q₀, R₀ of respective degrees s, d - s, such that P = Q₀R₀, \(\overline{\mathbf{Q}} = \overline{\mathbf{Q}}_0\) , \(\overline{\mathbf{R}} = \overline{\mathbf{R}}_0\) and Q₀ contains a monomial \(a_0X_1^{s}\) where φ(a) = φ(a₀). (Consider P, Q, R as polynomials in X₁ with coefficients in the ring B, the completion of A{[}X₂, \ldots, Xₙ{]} with respect to the valuation obtained by extending v by the method of Chapter VI, § 10, no. 1, Proposition 2; then apply Hensel's Lemma; finally use the initial hypothesis on P.)

\begin{enumerate}
\def\labelenumi{\arabic{enumi}.}
\setcounter{enumi}{15}
\item
  Let B be a discrete valuation ring whose residue field k is finite and is not a prime field; let \(k_{0}\) be the prime subfield of k and let A be the subring of B consisting of the elements whose classes in the residue field belong to \(k_{0}\) . Let \(\pi\) be a uniformizer of B and \((\theta_{i})_{1\leq i\leq m}\) a system of invertible elements of B such that the classes \(\bar{\theta}_{i}\) mod. \(\pi\) of the \(\theta_{i}\) form a system of representatives of \(k^{*}\) mod. \(k_{0}^{*}\) . Show that the elements \(p_{i}=\theta_{i}\pi\) and the element \(\pi\) are extremal in A and that every element in A is a product of an invertible element and powers of the \(p_{i}\) and \(\pi\) , although A is not integrally closed.
\item
  \begin{enumerate}
  \def\labelenumii{(\alph{enumii})}
  \tightlist
  \item
    Let A be an integral domain; show that in the ring \(A[X_{ij}]\) , where \((\mathbf{X}_{ij})\) is a family of \(n^{2}\) indeterminates \((1 \leqslant i \leqslant n, 1 \leqslant j \leqslant n)\) , the element \(\det(\mathbf{X}_{ij})\) is extremal. (Reduce it to the case where A is a field; observe that the factors of \(\det(\mathbf{X}_{ij})\) would necessarily be homogeneous polynomials and argue by induction on n.)
  \end{enumerate}
\end{enumerate}

\begin{enumerate}
\def\labelenumi{(\alph{enumi})}
\setcounter{enumi}{1}
\tightlist
\item
  Let \(\mathbf{K}\) be an infinite field and \(\mathbf{F}\) a polynomial in \(\mathbf{K}[\mathbf{Y}_1,\dots ,\mathbf{Y}_m]\) , which is also written as \(\mathrm{F(Y)}\) ; for every square matrix \(s = (\alpha_{ij})\) of order \(m\) with elements in K let \(F(s.Y)\) denote the polynomial F, where the element \(\sum_{j=1}^{m}\alpha_{ij}Y_{j}\) has been substituted for each \(Y_{i}\) . Show that, if F is extremal, so is \(F(s.Y)\) for every invertible matrix s. If there exists an integer \(k\geqslant0\) such that
\end{enumerate}

\[
\mathbf {F} (s. \mathbf {Y}) = (\det (s)) ^ {k} \mathbf {F} (\mathbf {Y})
\]

for every invertible matrix \(s\) , \(\mathbf{F}\) is necessarily homogeneous in each of the \(\mathbf{Y}_j\) ;

moreover, if \(\mathbf{F} = \mathbf{GH}\) where \(\mathbf{G}\) and \(\mathbf{H}\) are two polynomials in \(\mathbf{K}[\mathbf{Y}_1,\dots ,\mathbf{Y}_m]\) , there exist two integers \(p\) and \(q\) such that \(p + q = k\) and

\[
\mathbf {G} (s. \mathbf {Y}) = (\det (s)) ^ {p} \mathbf {G} (\mathbf {Y}), \quad \mathbf {H} (s. \mathbf {Y}) = (\det (s)) ^ {q} \mathbf {H} (\mathbf {Y})
\]

for every invertible matrix \(s\) (use (a)).

\begin{enumerate}
\def\labelenumi{(\alph{enumi})}
\setcounter{enumi}{2}
\tightlist
\item
  Let A be a ring which is not reduced to 0; consider the polynomial ring \(\mathbf{A}[\mathbf{X}_{ij}]\) , where the \(\mathbf{X}_{ij}\) are \(n(n + 1)/2\) (resp. \(2n(2n - 1)/2\) ) indeterminates with \(1 \leqslant i \leqslant j \leqslant n\) (resp. \(1 \leqslant i < j \leqslant 2n\) ); let \(U = (\xi_{ij})\) (resp. \(V = (\eta_{ij})\) ) be the square matrix of order \(n\) (resp. \(2n\) ) over \(\mathbf{A}[\mathbf{X}_{ij}]\) such that \(\xi_{ij} = \mathbf{X}_{ij}\) for \(1 \leqslant i \leqslant j \leqslant n\) and \(\xi_{ij} = \mathbf{X}_{ji}\) for \(i > j\) (resp. \(\eta_{ii} = 0\) for \(1 \leqslant i \leqslant 2n\) , \(\eta_{ij} = \mathbf{X}_{ij}\) for \(1 \leqslant i < j \leqslant 2n\) , \(\eta_{ij} = -\mathbf{X}_{ji}\) for \(i > j\) ). Show that \(\det(U)\) (resp. \(\operatorname{Pf}(V)\) ) is an extremal element in \(\mathbf{A}[\mathbf{X}_{ij}]\) (argue as in (b), considering \(\det(s, U, {}^ts)\) and \(\operatorname{Pf}(s, V, {}^ts)\) ).
\end{enumerate}

\begin{enumerate}
\def\labelenumi{\arabic{enumi}.}
\setcounter{enumi}{17}
\tightlist
\item
  Let K be a field and f = g/h an element of the field of rational functions \(\mathrm{K}(\mathrm{U}, \mathrm{V})\) in two indeterminates, where g and h are two relatively prime polynomials of \(\mathrm{K}[\mathrm{U}, \mathrm{V}]\) . Show that in the field of rational functions
\end{enumerate}

\[
\mathrm{K} \left(\mathrm{X} _ {1}, \mathrm{Y} _ {1}, \dots , \mathrm{X} _ {n}, \mathrm{Y} _ {n}\right)
\]

in \(2n\) indeterminates the determinant \(\operatorname{det}(f(\mathbf{X}_i, \mathbf{Y}_j))\) is equal to:

\[
\left(\prod_ {i, j} h \left(\mathrm{X} _ {i}, \mathrm{Y} _ {j}\right)\right) ^ {- 1} \mathrm{V} \left(\mathrm{X} _ {1}, \dots , \mathrm{X} _ {n}\right) \mathrm{V} \left(\mathrm{Y} _ {1}, \dots , \mathrm{Y} _ {n}\right) \mathrm{F} \left(\mathrm{X} _ {1}, \mathrm{Y} _ {1}, \dots , \mathrm{X} _ {n}, \mathrm{Y} _ {n}\right)
\]

where F is a polynomial in \(K[X_{1}, Y_{1}, \ldots, X_{n}, Y_{n}]\) and \(V(X_{1}, \ldots, X_{n})\) is the Vandermonde determinant (Algebra, Chapter III, § 6, no. 4). Consider the particular case where \(f = 1/(U + V)\) (``Cauchy's identity'').

\begin{enumerate}
\def\labelenumi{\arabic{enumi}.}
\setcounter{enumi}{18}
\tightlist
\item
  If \(U\) is a square matrix of order \(n\) , \(\Delta\) its determinant and \(\Delta_p\) the determinant of the \(p\) -th exterior power of \(U\) (Algebra, Chapter III, § 6, no. 3), show that:
\end{enumerate}

\[
\Delta_ {p} = \Delta^ {\binom {n - 1} {p - 1}}
\]

(use Exercise 11 of Algebra, Chapter III, § 6, and Exercise 17 (a) above).

\begin{enumerate}
\def\labelenumi{\arabic{enumi}.}
\setcounter{enumi}{19}
\tightlist
\item
  Let \(\mathbf{A}\) be a factorial domain and \(f = \sum_{k=0}^{n} a_k \mathbf{X}^k\) a polynomial in \(\mathbf{A}[\mathbf{X}]\) : suppose that there exists an extremal element \(p\) in \(\mathbf{A}\) such that:
\end{enumerate}

\begin{enumerate}
\def\labelenumi{(\arabic{enumi})}
\item
  there exists an index \(k \leqslant n\) such that \(a_{k}\) is not divisible by \(p\) but \(a_{i}\) is divisible by \(p\) for \(i < k\) ;
\item
  \(a_0\) is divisible by \(\pmb{p}\) but not by \(p^2\) .
\end{enumerate}

Show that under these conditions one of the irreducible factors of \(f\) in \(\mathbf{A}[\mathbf{X}]\) is of degree \(\geqslant k\) (argue in \((\mathbf{A} / p\mathbf{A})[\mathbf{X}]\) ). Consider the particular case where \(k = n\) (``Eisenstein's irreducibility criterion'').

\begin{enumerate}
\def\labelenumi{\arabic{enumi}.}
\setcounter{enumi}{20}
\tightlist
\item
  Show that in \(\mathbf{Z}[\mathbf{X}]\) the following polynomials are irreducible: \(\mathbf{X}^n - a\) , where one of the prime factors of \(a\) has exponent 1; \(\mathbf{X}^{2^k} + 1\) , (replace \(\mathbf{X}\) by \(\mathbf{X} + 1\) ); \(\mathbf{X}^4 + 3\mathbf{X}^3 + 3\mathbf{X}^2 - 5\) ; \(5\mathbf{X}^4 - 6\mathbf{X}^3 - a\mathbf{X}^2 - 4\mathbf{X} + 2\) . (Use Exercise 20.)
\end{enumerate}

¶ 22. (a) Let k be an ordered field and A the quotient of the polynomial ring k{[}X, Y, Z{]} by the principal ideal \((\mathbf{X}^{2} + \mathbf{Y}^{2} + \mathbf{Z}^{2} - 1)\) ; let x, y, z denote the canonical images of X, Y, Z in A. Show that A is a factorial domain (consider the ring of fractions \(A_{z-1}\) (notation of Chapter II, § 5, no. 1), using no. 4, Proposition 3).

\begin{enumerate}
\def\labelenumi{(\alph{enumi})}
\setcounter{enumi}{1}
\tightlist
\item
  Let \((e_i)_{1 \leq i \leq 3}\) be the canonical basis of \(\mathbf{A}^3\) and let \(\mathbf{M}\) be the quotient of \(\mathbf{A}^3\) by the monogenous sub-A-module \(\mathbf{N}\) generated by \(xe_1 + ye_2 + ze_3\) ; show that \(\mathbf{M}\) is a projective A-module (form a complementary submodule of \(\mathbf{N}\) in \(\mathbf{A}^3\) ). * (c) Show that, if \(k = \mathbf{R}\) , the A-module \(\mathbf{M}\) is not free (identify \(\mathbf{M}\) with a submodule of the module of continuous sections of the fibre bundle of tangent vectors to the unit sphere and use the fact that there exists no continuous field of tangent vectors \(\neq 0\) at every point of the sphere). *
\end{enumerate}

\begin{enumerate}
\def\labelenumi{\arabic{enumi}.}
\setcounter{enumi}{22}
\tightlist
\item
  Let B be a Krull domain, E its field of fractions, \(\Delta\) a derivation of E such that \(\Delta(\mathbf{B}) \subset \mathbf{B}\) , K the subfield of E the kernel of \(\Delta\) and A the Krull domain \(B \cap K\) ; suppose that K is of characteristic p \textgreater{} 0; then \(E^{p} \subset K\) and \(B^{p} \subset A\) , so that B is the integral closure of A in E and the canonical homomorphism \(\vec{i}: C(A) \to C(B)\) is defined ( \(\S 1\) , no. 10). Let U denote the group of invertible elements of B.
\end{enumerate}

\begin{enumerate}
\def\labelenumi{(\alph{enumi})}
\item
  If \(b \in \mathbf{E}\) is such that \(\mathrm{div}_{\mathbf{B}}(b)\) is the canonical image of a divisor in \(\mathbf{D}(\mathbf{A})\) , show that \(\Delta b / b \in \mathbf{B}\) (note that for every prime ideal \(\mathfrak{P}\) of \(\mathbf{B}\) of height 1 there exists \(b' \in \mathbf{K}\) such that \(v_{\mathfrak{P}}(b) = v_{\mathfrak{P}}(b')\) and observe that \(\mathbf{B}_{\mathfrak{P}}\) is invariant under \(\Delta\) ). Let \(\mathbf{L}\) be the additive subgroup of \(\mathbf{B}\) consisting of the \(\Delta b / b\) (``logarithmic derivatives'') which belong to \(\mathbf{B}\) (for \(b \in \mathbf{E}\) or \(b \in \mathbf{B}\) , which amounts to the same, and \(b \neq 0\) ) and let \(\mathbf{L}'\) be the subgroup of \(\mathbf{L}\) consisting of the \(\Delta u / u\) , where \(u \in \mathbf{U}\) . Show that, for every divisor \(d \in \mathbf{D}(\mathbf{A})\) such that the image of \(d\) is a principal divisor in \(\mathbf{D}(\mathbf{B})\) , the class mod. \(\mathbf{L}'\) of \(\Delta b / b\) for all \(b\) such that \(i(d) = \mathrm{div}_{\mathbf{B}}(b)\) depends only on the class of \(d\) in \(\mathbf{C}(\mathbf{A})\) and derive a canonical injective homophism \(\phi\) of \(\operatorname{Ker}(\bar{i})\) to \(\mathbf{L} / \mathbf{L}'\) .
\item
  Show that, if \(\Delta(B)\) is contained in no prime ideal of \(B\) of height 1 and \([E: K] = p\) , \(\phi\) is bijective. (Reduce it to showing that, if \(b \in E\) is such that \(\Delta b / b \in B\) and \(\mathfrak{P}\) is a prime ideal of \(B\) of height 1 such that \(v_{\mathfrak{P}}(b)\) is not a multiple of \(p\) , then \(e(\mathfrak{P} / \mathfrak{p}) = 1\) , where \(\mathfrak{p} = \mathfrak{P} \cap A\) ; for this, deduce from the hypothesis that, if \(t\) is a uniformizer of \(B_{\mathfrak{P}}\) , then \(\Delta t / t \in B_{\mathfrak{P}}\) , so that \(\mathfrak{P}B_{\mathfrak{P}}\) is invariant under \(\Delta\) and that \(\Delta\) therefore defines by taking quotients a derivation \(\overline{\Delta}\) of the residue field \(k = B_{\mathfrak{P}} / \mathfrak{P}B_{\mathfrak{P}}\) ; show that \(\overline{\Delta} \neq 0\) and deduce that \(f(\mathfrak{P} / \mathfrak{p}) = p\) .)
\item
  Let \(k\) be a field of characteristic 2, B the polynomial ring \(k[\mathbf{X}, \mathbf{Y}, \mathbf{Z}]\) and \(\Delta\) the derivation of \(\mathbf{E} = k(\mathbf{X}, \mathbf{Y}, \mathbf{Z})\) such that \(\Delta(\mathbf{X}) = \mathbf{Y}^4\) , \(\Delta(\mathbf{Y}) = \mathbf{X}^2\) , \(\Delta(\mathbf{Z}) = \mathbf{XYZ}\) ; the field \(\mathbf{K}\) the kernel of \(\Delta\) is such that \([\mathbf{E} : \mathbf{K}] = 4\) . Then \(\Delta(\mathbf{Z}) / \mathbf{Z} \in \mathbf{B}\) , but show that \(\mathrm{div}_{\mathbf{B}}(\mathbf{Z})\) is not the canonical image of a divisor in \(\mathrm{D}(\mathbf{A})\) . (Argue by reductio ad absurdum supposing that \(e(\mathfrak{P} / \mathfrak{p}) = 1\) for \(\mathfrak{P} = \mathbf{B}\mathbf{Z}\) , \(\mathfrak{p} = \mathfrak{P} \cap \mathbf{A}\) ; there would then be in \(\mathbf{A}\) a uniformizer of \(\mathrm{B}_{\mathfrak{P}}\) , necessarily of the form \(a(\mathbf{X}, \mathbf{Y}, \mathbf{Z})\mathbf{Z}\) where \(b = a(\mathbf{X}, \mathbf{Y}, 0) \neq 0\) ; deduce that \(\Delta b / b = -\mathbf{X}\mathbf{Y}\) and obtain a contradiction by calculating \(\Delta(-\mathbf{X}\mathbf{Y})\) .)
\end{enumerate}

\begin{enumerate}
\def\labelenumi{\arabic{enumi}.}
\setcounter{enumi}{23}
\tightlist
\item
  \begin{enumerate}
  \def\labelenumii{(\alph{enumii})}
  \tightlist
  \item
    Let E be a field of characteristic 2, \(\Delta\) a derivation of E and K the subfield of E the kernel of \(\Delta\) ; suppose that \([E:K]=2\) . Show that \(\Delta^{2}=a\Delta\) where \(a\in K\) and that for an element \(t\in E\) to be of the form \(\Delta x/x\) , it is necessary and sufficient that \(\Delta t=at+t^{2}\) .
  \end{enumerate}
\end{enumerate}

\begin{enumerate}
\def\labelenumi{(\alph{enumi})}
\setcounter{enumi}{1}
\tightlist
\item
  Let B be a factorial local integral domain of characteristic 2, m its maximal ideal, E its field of fractions and \(\Delta\) a derivation of E; suppose that the subfield K of E, the kernel of \(\Delta\) , satisfies {[}E: K{]} = 2; moreover, suppose that there exist two elements x, y of m such that \(\Delta x\) and \(\Delta y\) generate the ideal q of B generated by \(\Delta(B)\) . Show then that, if \(t = \Delta z / z\) belongs to q, there exists an invertible element u of B such that \(t = \Delta u / u\) (write \(t = r \Delta x + s \Delta y\) , where r, s are in B, and use (a)).
\end{enumerate}

¶ 25. Let k be a field of characteristic 2, B the ring of formal power series \(k[[X,Y]]\) , E its field of fractions and \(\Delta\) the k-derivation of E defined by \(\Delta(X)=Y^{2j}\) , \(\Delta(Y)=X^{2i}\) (i,j integers \(\geqslant0\) ); the subring A of B consisting of the \(x\in B\) such that \(\Delta x=0\) is the ring of formal power series in \(k[[X,Y,Z]]\) where \(X^{2}\) is substituted for X, \(Y^{2}\) for Y and \(X^{2i+1}+Y^{2j+1}\) for Z.

\begin{enumerate}
\def\labelenumi{(\alph{enumi})}
\item
  Show that the group C(A) contains a vector space over k of dimension \(N(i,j)\) equal to the number of ordered pairs of integers \((a,b)\) such that \(0 \leqslant a < i, 0 \leqslant b < j\) and \((2j + 1)a + (2i + 1)b \geqslant 2ij\) . (Use Exercise 23 (b) and Exercise 24 (a), note that the elements of L are the formal power series \(F \in B\) such that \(\Delta F = F^{2}\) ; attributing to X the weight \(2j + 1\) and to Y the weight \(2i + 1\) , decompose F as an infinite sum of isobaric polynomials; if \(L_{q}\) is the subgroup of L consisting of the \(F \in L\) whose isobaric components are of weight \(\geqslant q\) , \(L/L'\) is isomorphic to the direct sum of the groups \(C_{q}/C_{q+1}\) , where \(C_{q} = L_{q}/(L' \cap L_{q})\) ; calculate these groups for \(q \geqslant 4ij\) , using Exercise 24 (b).)
\item
  Show that the ideal \(AX^{2}\) of A is prime; if \(A' = A[X^{-2}]\) , deduce that C(A') and C(A) are isomorphic (no. 4, Proposition 3). Show that A' is a Dedekind domain (consider the ring \(B' = B[X^{-2}]\) , which is integral over A' and a principal ideal domain, and use Chapter V, § 2, no. 4, Theorem 3). Deduce an example of a Dedekind domain whose ideal class group is infinite.
\end{enumerate}

\begin{enumerate}
\def\labelenumi{\arabic{enumi}.}
\setcounter{enumi}{25}
\tightlist
\item
  Let A be a factorial domain and K its field of fractions. For elements \(i\) ( \(1 \leqslant i \leqslant r\) ) of B = A{[}X \(_{1}\) , \ldots, X \(_{n}\) {]} to be such that the ideal \(\sum_{i=1}^{r} \text{Bf}_i\) is equal to B, it is necessary and sufficient that there exist polynomials \(v_{i}\) ( \(1 \leqslant i \leqslant r\) ) in \(\mathbf{K}[\mathbf{X}_1, \ldots, \mathbf{X}_n]\) such that \(\sum_{i=1}^{r} v_i f_i = 1\) and which satisfy further the following condition: if we write \(v_i = w_i / d\) , where \(d \in \mathbf{A}\) , the polynomials \(w_i\) belong to \(\mathbf{A}[\mathbf{X}_1, \ldots, \mathbf{X}_n]\) and the g.c.d. of the set of coefficients of all the \(w_i\) ( \(1 \leqslant i \leqslant r\) ) is equal to 1, then, for every extremal element \(p\) of A dividing \(d\) , the ideal generated by the classes of the \(f_i\) in the ring \((\mathrm{A} / \mathrm{A}p)[\mathrm{X}_1, \ldots, \mathrm{X}_n]\) is the whole of this ring.
\end{enumerate}

¶ 27. (a) Let K be a field and B the ring generated, in an algebraic closure of the field of rational functions K(U, V, X, Y) in 4 indeterminates, by the polynomial ring K{[}U, V, X, Y{]} and a root z of the polynomial

\[
\mathrm{F} = \mathrm{Z} ^ {7} - \mathrm{U} ^ {5} \mathrm{X} ^ {2} - \mathrm{V} ^ {4} \mathrm{Y} ^ {3}.
\]

Show that \(\mathbf{B}\) is a factorial domain (cf.~Exercise 7).

\begin{enumerate}
\def\labelenumi{(\alph{enumi})}
\setcounter{enumi}{1}
\tightlist
\item
  Let p be the (prime) ideal generated by X, Y, U, V and z in A and write \(C = A_{p}[[T]]\) ; show that C is a Noetherian local integral domain, which is not factorial, but whose associated graded domain \(\text{gr}(C)\) is factorial (use Exercise 8).
\end{enumerate}
% semantic-fragments: legacy-input-seam
\relax{}
